% ============================================================
\section{Conclusion}
\label{sec:conclusion}
% ============================================================

Representation drift is widely read as forgetting.  Against the intervention it
stands in for, on three time-series foundation models and six benchmarks, that
reading does not hold: the least drifted cells can be the ones that degrade, a
representation driven nearly orthogonal can retain its capability, and CKA adds
nothing to what backbone and protocol already predict.  On the target task we
could not find one cell where fine-tuning demonstrably damaged a pre-trained
capability; on datasets it was not tuned on we found eleven, and CKA did not order
them.  None of this makes CKA a poor
measure of similarity; it makes similarity a poor measure of harm.

The practical advice follows from the two measurement choices that decided our
results.  First check that the pre-trained model beats a baseline fitted on your
own data, because a forgetting result on a cell with no pre-trained advantage is a
statement about the benchmark.  Then compare full fine-tuning against a strictly
frozen encoder on windows neither run selected on.  That comparison costs one
extra training run, and it answers the question drift only appears to.
